\section{Arithmetic Euler characteristic and finiteness for the actual cubic frame}
\label{sec:cubic-arithmetic-euler-phi}

The Euler-characteristic field of the source frame quantifies over
every finite coefficient field, every finite-dimensional continuous
representation, and every invariant local subspace.  The argument
below verifies that full statement for the actual away-inertia
quotient.  It does not require the coefficient representation to be
split by the fixed residual splitting field.

\paragraph{Arithmetic group and full local maps.}
Let \(p>2\), let \(F\) be a totally real number field, and let
\(L/F\) be a finite Galois extension of degree prime to \(p\)
containing \(\mu_p\).  In the application, take \(L\) to be the
field cut out by the full residual tuple, the CM quadratic
character and \(\mu_p\).  Inside a fixed
absolute Galois group \(G_F\), let
\[
 N=\overline{\langle\!\langle
       I_w:\ w\text{ a finite place of }L,\ w\nmid p
       \rangle\!\rangle},\qquad
 \Gamma=G_F/N,\qquad \Delta=\operatorname{Gal}(L/F).
 \tag{G96.1}
\]
All conjugate inertia groups are included in this definition.
Their closed normal closure inside \(G_L\) is stable under \(G_F\),
because conjugation permutes that entire collection.  Thus it is
also their closed normal closure inside \(G_F\), and
\[
 1\longrightarrow G_{L,S_p}\longrightarrow\Gamma
   \longrightarrow\Delta\longrightarrow1.
 \tag{G96.2}
\]
Here \(G_{L,S_p}\) is the Galois group of the maximal extension
of \(L\) unramified outside the places above \(p\) and infinity.
For that actual choice of \(L\), (G96.1) kills every away-\(p\)
inertia intersection with the kernel of the full finite residual
datum: adjoining \(\mu_p\) does not change these intersections,
since its inertia away from \(p\) is trivial.  This is the
quotient identified in AC-G90-quotient.

For every \(v\mid p\), retain the full absolute local group
\(G_{F_v}\) and its continuous map
\(\iota_v:G_{F_v}\to\Gamma\).  We do not replace it by its image
in \(\Gamma\).  Put \(d_v=[F_v:\mathbb Q_p]\), so
\(\sum_{v\mid p}d_v=[F:\mathbb Q]\), and let \(c_w\in\Gamma\)
be the image of complex conjugation for each real embedding
\(w\) of \(F\).  The cyclotomic character, and its reduction
\(\bar\epsilon\), factor through \(\Gamma\), since they are
unramified away from \(p\).

\paragraph{The literal cone.}
Let \(k\) be any finite field of characteristic \(p\), let \(W\)
be a finite-dimensional continuous \(k[\Gamma]\)-module, and let
\(b_v\subset W\) be a \(G_{F_v}\)-stable \(k\)-subspace.  Write
\[
 C_b=\operatorname{Cone}\left(
 C^\bullet_{\rm cont}(\Gamma,W)\longrightarrow
 \bigoplus_{v\mid p}C^\bullet_{\rm cont}(G_{F_v},W/b_v)
 \right)[-1].
 \tag{G96.3}
\]
All coefficients in these complexes have the discrete topology.
The finite target makes continuous cochains locally constant.
In degrees zero, one and two, (G96.3) gives exactly the source
\path{coneH0}, \path{H1C}, and \path{coneH2} in
\path{ContinuousCoboundaries2.lean}.  In degree one a cocycle is
\((z,(t_v\bmod b_v))\) with
\(z(\iota_v s)-\{s t_v-t_v\}\in b_v\), modulo
\((dX,(X\bmod b_v))\).  In degree two the condition is
\(\operatorname{res}_v c-du_v\in b_v\), and boundaries have
\(c=d\beta\), \(u_v=\operatorname{res}_v\beta-dY_v\bmod b_v\).
These are the displayed source formulas, with no change to the
local condition or to its degree-zero primitives.  Conversely,
cochains in \(W/b_v\) lift continuously to \(W\), by a fixed
linear section of this finite quotient, proving the identifications
in both directions.

\begin{theorem}[AC-G96-EC: arithmetic verification of the full source EC field]
\label{AC-G96-EC}
In (G96.1)--(G96.3), suppose
\[
 \operatorname{Hom}_{k[\Gamma]}(W,k(\bar\epsilon))=0.
 \tag{G96.4}
\]
Then all the cohomology groups of \(C_b\) are finite-dimensional,
they vanish outside degrees zero through two, and
\[
 \begin{split}
 \dim_k H^2(C_b)+\dim_k H^0(C_b)
       +\sum_{v\mid p}d_v\dim_k b_v
  =\dim_k H^1(C_b)+\sum_{w\mid\infty}\dim_k W^{c_w}.
 \end{split}
 \tag{G96.5}
\]
Consequently the arithmetic frame with these full local maps
satisfies the literal source \path{Frame.EC}, for every coefficient
field and representation in that definition.  Its asserted
inequality holds with equality.  This conclusion requires neither
an abelian host nor an estimate for any class-group sector.
\end{theorem}

\begin{proof}
\emph{1. Realize the actual quotient by a constructible sheaf.}
Set \(X=\operatorname{Spec}\mathcal O_F[1/p]\) and
\(Y=\operatorname{Spec}\mathcal O_L[1/p]\), and let
\(\pi:Y\to X\) be the finite normalization map.  Restrict \(W\)
to the normal subgroup \(G_{L,S_p}\) in (G96.2).  The resulting
representation defines a locally constant finite sheaf
\(\mathcal W_L\) on \(Y\), with the \(\Delta\)-equivariance
induced by the original \(\Gamma\)-action.  This construction
does not assert that \(G_{L,S_p}\) acts trivially on \(W\).
Put
\[
       \mathcal F_W=(\pi_*\mathcal W_L)^\Delta.
 \tag{G96.6}
\]
Finite pushforward on the small etale site is exact, has no
higher direct images, and has at a geometric point the direct sum
of the stalks above that point.  These are Stacks Project,
Tags \texttt{03QO} and \texttt{03QP}: after strict henselization
the finite source is a finite disjoint union of strictly henselian
local schemes, whose global-section functor is the closed stalk.
Taking \(\Delta\)-invariants is the image of the idempotent
\(|\Delta|^{-1}\sum_{\delta\in\Delta}\delta\), hence is exact.
The arithmetic sheaf/group comparison in Milne,
\emph{Arithmetic Duality Theorems}, II, Proposition 2.9, and
Hochschild--Serre for (G96.2) therefore give natural comparisons
\[
 \begin{split}
 R\Gamma(X,\mathcal F_W)
   &\simeq R\Gamma(Y,\mathcal W_L)^\Delta\simeq
     R\Gamma(\Gamma,W),\\
 H^i(X,\mathcal F_W)
   &\simeq H^i(G_{L,S_p},W)^\Delta.
 \end{split}
 \tag{G96.7}
\]
The complexes can be compared using the same bar resolutions;
the positive-degree group cohomology of \(\Delta\) vanishes by
averaging.  In particular this is a derived comparison, not a
dimension division by \(|\Delta|\).

Choose an open immersion \(j:U\hookrightarrow X\) omitting the
finitely many away-\(p\) primes ramified in \(L/F\).
On \(U\), \(W\) is a locally constant sheaf.  The geometric stalk
of (G96.6) at an omitted prime is \(W^{I_v}\): the finite
pushforward stalk is the permutation sum over the geometric
places above \(v\), and its invariant part is the stabilizer
invariant part.  The stabilizer action is exactly the finite
inertia action on \(W\), since its intersection with \(G_L\)
is killed in \(\Gamma\).  Consequently
\[
                   \mathcal F_W=j_*(W|_U).
 \tag{G96.8}
\]
This also proves constructibility.  The natural comparisons
(G96.7) commute with the maps to the full local groups at \(p\).
One can check this after restriction to the places of \(L\),
where it is the usual natural arithmetic comparison; restriction
to their full sum is injective on the relevant invariant groups,
because the local indices divide \(|\Delta|\).

\emph{2. The strict cone and its top cohomology.}
Write \(C_0\) for (G96.3) with every \(b_v=0\).
Milne's compact-support definition and II, Proposition 2.3
identify it, using (G96.7), with
\[
                         C_0\simeq R\Gamma_c(X,\mathcal F_W).
 \tag{G96.9}
\]
At real places this definition uses Tate cochains.  Their
\(p\)-primary complexes are acyclic because \(p>2\), so they
add no local terms to our full \(p\)-local cone.
All its cohomology groups are finite by II, Theorem 2.13.
Moreover (G96.2), exact invariants, and the global
cohomological-dimension statement in I, Theorem 4.10(c), give
\(H^i(\Gamma,W)=0\) for \(i>2\).  The local absolute groups
have \(p\)-cohomological dimension two.  Thus \(C_0\) is
concentrated cohomologically in degrees zero through three.

Apply Artin--Verdier duality in the precise form II,
Theorem 3.1(b), with \(r=0\), to the constructible sheaf
\(\mathcal F_W\).  It gives a perfect finite pairing
\[
 \operatorname{Hom}_X(\mathcal F_W,\mathbb G_m)
       \times H^3_c(X,\mathcal F_W)\longrightarrow\mathbb Q/\mathbb Z.
 \tag{G96.10}
\]
Every map in the first factor has image in \(\mu_p\).
Restriction to \(U\) identifies that factor with
\(\operatorname{Hom}_{\mathbb F_p[\Gamma]}(W,\mathbb F_p(1))\).
Indeed a generic map extends by applying \(j_*\), since
\(j_*\mu_p=\mu_p\) on the normal scheme \(X\); conversely a map
to the locally constant sheaf \(\mu_p\) is determined on the
dense open \(U\).  Formula (G96.8) makes both assertions literal.
The finite-field trace pairing identifies
\[
 \operatorname{Hom}_{k[\Gamma]}(W,k(\bar\epsilon))
   \ \simeq\
 \operatorname{Hom}_{\mathbb F_p[\Gamma]}(W,\mathbb F_p(1))
 \tag{G96.11}
\]
as \(\mathbb F_p\)-spaces: send \(f\) to
\(\operatorname{Tr}_{k/\mathbb F_p}\circ f\); the inverse is
characterized by
\(\operatorname{Tr}(a f(w))=\ell(aw)\) for all \(a\in k\).
Thus (G96.4) implies \(H^3(C_0)=0\).
Only the stated duality theorem is imported here; we do not
claim a new proof of Artin--Verdier duality.

\emph{3. Arbitrary stable local subspaces.}
Continuous cochains in a finite quotient lift continuously,
so there is an exact sequence of complexes
\[
 0\longrightarrow
       \bigoplus_{v\mid p}C^\bullet(G_{F_v},b_v)[-1]
 \longrightarrow C_0\longrightarrow C_b\longrightarrow0.
 \tag{G96.12}
\]
Its long exact sequence makes \(H^3(C_b)\) a quotient of
\(H^3(C_0)\), since \(H^3(G_{F_v},b_v)=0\).  Hence
\(H^3(C_b)=0\).  All other groups are finite, and vanish in
degrees below zero and above three, directly from (G96.3).
This proves the asserted range and finiteness for \(C_b\).

\emph{4. Calculate the exact Euler characteristic.}
For a finite complex use the additive \(k\)-dimension Euler
characteristic.  Milne II, Theorem 2.13(b), evaluated at the real
places with the just specified Tate convention, gives
\[
             \chi(C_0)=\sum_{w\mid\infty}\dim_k W^{c_w}.
 \tag{G96.13}
\]
For clarity, the ordinary global version is
\(\chi(\Gamma,W)=\sum_w\dim_k W^{c_w}-[F:\mathbb Q]\dim_k W\),
and subtracting the full \(p\)-local Euler characteristics gives
the same (G96.13).  This also follows from II, Theorem 2.13(a)
or its proof by removing the finitely many ramified closed points.
Each removed finite-field term and each away-\(p\) local term
has Euler characteristic zero.

The local Euler formula, Milne I, Theorem 2.8, gives
\(\chi(G_{F_v},b_v)=-d_v\dim_k b_v\).
Taking Euler characteristics in (G96.12) therefore yields
\[
 \chi(C_b)=\chi(C_0)+\sum_{v\mid p}\chi(G_{F_v},b_v)
   =\sum_w\dim_k W^{c_w}-\sum_{v\mid p}d_v\dim_k b_v.
 \tag{G96.14}
\]
There is no degree-three term by Step 3.  Expanding the left
side as \(h^0-h^1+h^2\) gives (G96.5).
Finally, the source's open-kernel condition makes the action on
the finite set \(W\) continuous, its local stability condition is
exactly the one used here, and its condition on linear maps is
(G96.4).  Its last sum is \(\sum_w\dim_k\ker(\rho(c_w)-1)\).
This verifies every quantified field of \path{Frame.EC}.
\end{proof}

\begin{lemma}[Arithmetic finiteness: a direct verification of the AC-G94-space input]
\label{G96-phi-verification}
The group \(\Gamma\) in (G96.1), and every open subgroup of it,
satisfy the literal source \path{PhiCond} at \(p\):
\[
     \#\operatorname{Hom}_{\rm cont}(U,\mathbb F_p)<\infty
       \quad\text{for every open }U\subseteq\Gamma.
 \tag{G96.15}
\]
In particular this supplies the finiteness hypothesis for the
changed-residual universal ring of AC-G94-space on the CM
index-two subgroup, as well as the \path{Frame.Ok.phi} field.
\end{lemma}

\begin{proof}
Take a finite set \(S\) of primes of \(F\) containing the primes
above \(p\), infinity, and all primes ramified in \(L/F\).
At any finite \(v\notin S\), full inertia lies in \(G_L\) and is
killed by (G96.1).  Hence the quotient map factors as a continuous
surjection \(G_{F,S}\twoheadrightarrow\Gamma\).
For an open subgroup \(U\), its preimage is the Galois group
\(G_{F',S'}\) for a finite extension \(F'/F\) in the maximal
extension unramified outside \(S\), with \(S'\) the places above
\(S\).  The usual maximality statement holds also when \(F'/F\)
is not Galois: the normal closure over \(F\) of an extension
unramified outside \(S'\) is still unramified outside \(S\).
Precomposition gives an injection
\[
 \operatorname{Hom}_{\rm cont}(U,\mathbb F_p)
       \hookrightarrow H^1(G_{F',S'},\mathbb F_p).
 \tag{G96.16}
\]

Here is the finite arithmetic calculation on the right.  Put
\(E=F'(\mu_p)\), of degree prime to \(p\) over \(F'\), and
let \(S_E\) be the places above \(S'\).  Restriction on this
\(H^1\) is injective: \(\operatorname{cor}\circ\operatorname{res}
=[E:F']\) on the original cohomology, and this degree is
prime to \(p\).  Over \(E\) identify
\(\mathbb F_p\) with \(\mu_p\).  The Kummer sequence on
\(\operatorname{Spec}\mathcal O_{E,S_E}\), together with
Milne II, Proposition 2.9, gives
\[
 0\longrightarrow \mathcal O_{E,S_E}^{\times}/
                       (\mathcal O_{E,S_E}^{\times})^p
 \longrightarrow H^1(G_{E,S_E},\mu_p)
 \longrightarrow \operatorname{Cl}(\mathcal O_{E,S_E})[p]
 \longrightarrow0.
 \tag{G96.17}
\]
The first group is finite by finite generation of the group
of \(S_E\)-units; the last is finite by class-group finiteness.
Thus (G96.16) is finite.  These are qualitative finiteness
inputs, not class-number bounds.

A homomorphism from a profinite group to the finite discrete
group \(\mathbb F_p\) is continuous if and only if its kernel
is open.  Rewriting the additive target as the source's
\path{Multiplicative (ZMod p)} therefore gives exactly
\path{PhiCond}.  An open subgroup of an open subgroup is open,
so the assertion persists on the CM subgroup and on every finite
open arithmetic layer.  This is also the quotient/permutation
cohomology mechanism of the inspected source proofs
\path{phi_condition_of_finite_cohomology} and \path{phi_surj}
in \path{RingEx.lean}; no residual irreducibility is used.
\end{proof}

\paragraph{Application and limits.}
For the selected-place genus host and its cyclotomic layers in
AC-G93-fill, take at each real field its finite prime-to-\(p\)
Galois field splitting the full residual tuple, the CM character,
and \(\mu_p\).  The same construction gives (G96.1), with the
full local groups at all coefficient-prime places.  Neither proof
requires these completions to be unramified over \(\mathbb Q_p\).
Thus \path{EC} and \path{PhiCond} are verified on those actual
groups and on the changed-residual space.  The upper bounds
\path{b1p}, \path{b1}, and \path{h1}, the inequalities in
\path{Big}, the remaining source \path{Frame.Ok} fields, and
the per-curve Hecke-ideal containment are separate assertions.

\paragraph{Geometric dictionary and source scope.}
The finite residual field gives a finite normalization of the
arithmetic curve; its finite local systems and exact prime-to-\(p\)
invariants realize the prescribed away-inertia quotient.  The
source local cone is compact-support cohomology followed by
allowing the specified local subspaces.  Stacks Tags
\texttt{03QO}/\texttt{03QP}, Milne I.2.8 and II.2.3/2.9/2.13,
and the exact statement of II.3.1(b) are the applied inputs.
The finite-cover, local-system and restriction interpretation
agrees with the already audited SGA~1 category; no additional
Grothendieck theorem or new proof of arithmetic duality is
attributed here.  No Lean execution or formalization of these
new arithmetic arguments is claimed.
